\subsection{Invertibility and Isomorphisms}

I think it was nice that change of basis was presented before eigenvectors\dots

The notes for this chapter were pretty long but all the exercises were pretty easy, with none of them that I attempted needing a google or anything.

\begin{enumerate}
    \item[1] Suppose $T \in \mathcal L(V, W)$ is invertible. Show that $T^{-1}$ is invertible and that this inverse is $(T^{-1})^{-1}$.
    
    \begin{proof}
        Well, we would want 

        \[T^{-1}(T^{-1})^{-1} = I\]

        and

        \[(T^{-1})^{-1}T^{-1}=I\]

        and $T$ satisfies both, so it is the inverse of $T^{-1}$.
    \end{proof}

    \item[3] Suppose $V$ is finite-dimensional and $T \in \mathcal L(V)$. Prove that the following are equivalent.
    
    \begin{enumerate}
        \item T is invertible
        \item $Tv_1, \dots Tv_n$ is a basis of V for every basis $v_1, \dots, v_n$.
        \item $Tv_1, \dots, Tv_n$ is a basis for $V$ for some basis $v_1, \dots, v_n$ of $V$.
    \end{enumerate}

    \begin{proof}
        We prove (a) and (b) are equivalent by showing that they are both equivalent to injectivity. (a) is already proven to be equivalent to injectivity. 

        Suppose that $Tv_1, \dots, Tv_n$ is a basis for $V$ for some basis $v_1, \dots, v_n$ of $V$. Suppose we have $u, v, w \in V$ where $Tv = w$ and $Tu = w$. $T(v - u) = Tv - Tu = 0$. If $v \neq u$, then that implies that we could construct a non-zero linear combination of $Tv_1, \dots, Tv_n$ that equaled 0, which is a contradiction. 

        Suppose that $T$ is injective. This implies surjectivity, so each $w \in V$ has some unique $u \in V$ such that $Tu = w$. Given some basis of $V$, $v_1, \dots, v_n$, $Tv_i$ forms a basis since there is a unique linear combination for all elements of $V$, the definition of a basis.

        Notice I did not mention any particular basis, which means it works for any and all bases.
    \end{proof}
    

    \item[5] Suppose $V$ is finite-dimenstional, $U$ is a subspace of $V$, and $S \in \mathcal L(U, V)$. 
    
    Prove that there exists an invertible linear map $T$ from $V$ to itself such that $Tu = Su$ for every $u \in U$ if and only if $S$ is injective.

    \begin{proof}
        Suppose $S$ is not injective. Then $T$ cannot be injective either, and so it also cannot be invertible.

        Suppose that $S$ is injective. Let $u_1, \dots, u_m$ be a basis for $U$. Then, $Su_1, \dots Su_m$ is a basis for some subspace of V. Both can be extended and we can construct some basis for $V$, $u_1, \dots, u_m, v_1, \dots, v_n$ where the list $Tu_1, \dots Tu_m, Tv_1, \dots, Tv_n$ are bases for $V$. Then, $T$ will be invertible, so we have constructed a solution.
    \end{proof}


    \item[6] Suppose that $W$ is finite-dimensional and $S, T \in \mathcal L(V, W)$.
    
    Prove that $\lnull S = \lnull T$ if and only if there exists an invertible $E \in \mathcal L(W)$ such that $S = ET$.

    \begin{proof}
        Suppose there exists an invertible $E$ such that $S = ET$. For $v \in V$, if $T$ maps $v$ to 0, then so must $S$ because $E(0) = 0$. If $S$ maps $v$ to 0 then so must $T$ because $E(v) = 0$ if and only if $v = 0$ (because it is injective).


        Suppose that $\lnull S = \lnull T$. Consider the basis of $\lnull S$ (or $\lnull T$, they are the same), $v_1, \dots, v_n$. We can extend this basis to all of $V$ as $v_1, \dots, v_n, u_1, \dots, u_m$. Let the subspace formed from the basis $u_1, \dots, u_m$ be $U$. 
        
        We will show that $S$ and $T$ are injective in $U$. Suppose that there is some $u_1, u_2$ such that $Su_1 = Su_2$. Then, we would have $S(u_1 - u_2) = 0$, but $u_1 - u_2$ is in $\lnull S$. The only element in both $U$ and $\lnull S$ is 0, therefore $u_1 - u_2 = 0$ and $u_1 = u_2$. Same logic applies to $T$.
        
        Since $T$ is injective in $U$, $Tu_1, \dots, Tu_m$ are linearly independent in $W$. We can arbitrarily extend this to a basis $Tu_1, \dots, Tu_m, w_1, \dots w_r$, where $\dim W = m + r$. We also define a basis $Su_1, \dots, Su_m, p_1, \dots, p_r$ using the same principle.
        
        Now we can define $E$. 

        \[E(Tu_i) = Su_i\]

        and 

        \[E(w_i) = p_i\]

        $E$ is well defined because $Tu_1, \dots, Tu_m, w_1, \dots, w_r$ is a basis in $W$. 
        
        $E$ is invertible because $Su_1, \dots, Su_m, p_1, \dots, p_r$ is a basis in $W$. $S = ET$ because the bases all map to the same values (no I will not show it but it is obvious).
    \end{proof}

    \item[9] Suppose $V$ is finite-dimensional and $T: V \to W$ is a surjective linear map of $V$ onto $W$. Prove that there is a subspace $U$ of $V$ such that $T_{|U}$ is an isomorphism of $U$ onto $W$. 
    
    \quad \emph{Here $T_{|U}$ means that the function $T$ is restricted to $U$. Thus $T_{|U}$ is the function whose domain is $U$, with $T_{|U}$ defined by $T_{|U}(u) = Tu$ for every $u \in U$.}

    \begin{proof}
        Consider arbitrary basis of $W$, $w_1, \dots, w_m$. Because $T$ is surjective, we can find some vectors $v_1, \dots, v_m$ such that $Tv_i = w_i$ for all $i \in \{1, \dots, m\}$. $v_1, \dots, v_m$ is linearly independent, since if it wasn't, $w_1, \dots, w_m$ wouldn't be a basis. Now we define $U$ as the subspace formed from $v_1, \dots, v_m$ as a basis. $\dim \lrange T = \dim W$, and it remains $\dim W$ restricted to $U$ since $w_1, \dots, w_m$ is basis for $W$.


        \begin{align*}
            \dim U &= \dim \lrange T_{|U} + \dim \lnull T_{|U} \\
            m &= m + \dim \lnull T_{|U} \\
            \dim \lnull T_{|U} &= 0
        \end{align*}

        Therefore $T_{|U}$ is an invertible linear map, the definition of isomorphism.
        
    \end{proof}

    \item[12] Suppose $V$ is finite-dimensional and $S, T, U \in \mathcal L(V)$ and $STU = I$. Show that $T$ is invertible and that $T^{-1} = US$.
    
    \begin{proof}
        If $T$ mapped some non-zero vector to 0, there would be no way to recover it back to a non-zero value, so $T$ is invertible. This also shows that $U$ and $S$ are invertible.
        
        The inverse of $U$ is $ST$ and the inverse of $S$ is $TU$. 

        \begin{align*}
            STU = I \\
            STU(ST) = ST \\
            ST(US)T = ST \\
            S^{-1}T(US)TT^{-1} = S^{-1}STT^{-1} \\
            T(US) = I \\
            US = T^{-1}
        \end{align*}

        Therefore we done.
    \end{proof}


    \item[15] Suppose $T \in \mathcal(V)$ and $v_1, \dots, v_m$ is a list in $V$ such that $Tv_1, \dots, Tv_m$ spans $V$. Prove that $v_1, \dots, v_m$ spans $V$. 
    
    \begin{proof}
        Suppose that $v_1, \dots, v_m$ was not linearly independent. Then, there is some $a_1v_1+ \dots+ a_mv_m = 0$ with some $a_i \neq 0$. Then, $T(a_1v_1 + \dots + a_mv_m) = a_1Tv_1 + \dots a_mTv_m = 0$, which is a contradiction since $Tv_1, \dots, Tv_m$ is supposed to be a basis. Thus $v_1, \dots, v_m$ is linearly independent. By the linear map lemma, it also spans $V$ so it is a basis.
    \end{proof}


    \item[18] Show that $V$ and $\mathcal L(\F, V)$ are isomorphic vector spaces.
    
    \begin{proof}
        Well, $\F$ is just $\F^1$, and any map $T \in \mathcal L(\F, V)$ is entirely defined by $T1$, since all other numbers are scalar multiples of $T1$. So we can define our function $S \in \mathcal L(\mathcal L(\F, V), V)$ as $S(T) = T1$. This function has a clear inverse (that being taking some vector $v$ to the function that maps $1$ to $v$), so we have an isomorphism.
    \end{proof}

    \item[20] Suppose $q \in \mathcal P(\R)$. Prove that there exists a polynomial $p \in \P(\R)$ such that $q(x) = (x^2 + x)p''(x) + 2xp'(x) + p(3)$ for all $x \in \R$.
    
    \begin{proof}
        Define $T \in \P(\R)$ where $Tp = (x^2 + x)p'' + 2xp' + p(3)$. We want to prove that $T$ is surjective in $\P(\R)$.

        Consider the basis $1, x, \dots, x^n$ for vector space $\P_n(\R)$. $Tx^n$ is linearly independent from all previous polynomials because it is the only polynomial of degree $n$. Thus, $T1, Tx, \dots, Tx^n$ forms a basis, so it is surjective, so we can always construct $q$.
    \end{proof}

    \item[24] Suppose $A$ and $B$ are square matrices of the same size and $AB = I$. Prove that $BA = I$.
    
    \begin{proof}
        \begin{align*}
            AB = I \\
            B(AB) = B \\
            (BA)B = B \\
            BA = I 
        \end{align*}
    \end{proof}
\end{enumerate}